\documentclass[11pt]{article}
\usepackage[margin=1.05in]{geometry}
\usepackage{amsmath,amssymb,amsthm,mathtools}
\usepackage{booktabs}
\usepackage{enumitem}
\usepackage{tikz-cd}
\usepackage[hidelinks]{hyperref}

\newtheorem{theorem}{Theorem}
\newtheorem{lemma}{Lemma}
\newtheorem{proposition}{Proposition}
\newtheorem{corollary}{Corollary}
\theoremstyle{definition}
\newtheorem{definition}{Definition}
\newtheorem{construction}{Construction}
\theoremstyle{remark}
\newtheorem*{remark}{Remark}
\newtheorem*{convention}{Convention}

\newcommand{\Hs}{\mathcal H}
\newcommand{\K}{\mathcal K}
\newcommand{\B}{\mathcal B}
\newcommand{\Ecat}{\mathcal E}
\newcommand{\Oc}{\mathbb O}
\newcommand{\Sw}{\mathrm{Sw}}
\newcommand{\Ctx}{\mathsf{Ctx}}
\newcommand{\Gpd}{\mathbf{Gpd}}
\newcommand{\Cat}{\mathbf{Cat}}
\newcommand{\U}{\mathrm{U}}
\newcommand{\PU}{\mathrm{PU}}
\newcommand{\SU}{\mathrm{SU}}
\newcommand{\Wop}{W}
\newcommand{\id}{\mathrm{id}}
\newcommand{\tr}{\operatorname{tr}}
\newcommand{\hol}{\mathrm{hol}}
\newcommand{\Zt}{\mathbb Z_2}
\newcommand{\pushf}[1]{{#1}_{!}}

% status tags
\newcommand{\STC}{{\normalfont\textsc{[Complete]}}}
\newcommand{\STS}{{\normalfont\textsc{[Sketch]}}}
\newcommand{\STG}{{\normalfont\textsc{[Gap]}}}

\title{The Observer--Context Category:\\ Formal Architecture}
\author{Manuel Flores Gordillo}
\date{July 2026 --- standalone architecture note (companion to Paper D)}

\begin{document}
\maketitle

\begin{abstract}
Paper D introduces an ``observer'' as an isometric enlargement of a system together
with a measurement context on the enlarged algebra, with two arrow classes:
invertible \emph{switches} (context changes at a fixed cut) and non-invertible
\emph{refinements} (moving the Heisenberg cut, $A\mapsto A\otimes I$). The
non-invertibility of refinements is what forced the correction from ``groupoid'' to
``category.'' This note supplies the missing architecture. We argue that the correct
home is a \emph{Grothendieck opfibration} $\pi:\Oc\to\Ecat$ over the category
$\Ecat$ of enlargements --- equivalently a pseudofunctor $\Ctx:\Ecat\to\Gpd$ whose
fibers are the switch groupoids and whose reindexing functors are the
adjoin-and-tensor maps. This choice (i) subsumes the ``category with a wide switch
subgroupoid'' picture (the subgroupoid is exactly the union of fibers), (ii) is more
faithful than a bare double category because it builds the invertibility asymmetry
into the variance rather than positing it, and (iii) makes ``the cocycle is pulled
along a functor'' literal. We specify objects, all arrow types with sources and
targets, composition, identities, and associativity; prove the switches form a wide
subgroupoid; give refinement functoriality on contexts, on the $\U(1)$ cocycle, on
value assignments, and on obstruction classes with pullback--monotonicity; locate the
$\U(1)$ cocycle precisely (fiberwise, on the switch subgroupoid, as a
reindexing-compatible section); define holonomy only for invertible switch loops and
separate it from pullback-along-embeddings; and restate the movable-cut result as an
exact five-level statement with a commutative diagram. Every claim is labeled
\STC{} complete, \STS{} a sketch relying on a named categorical result, or \STG{} a
gap. The Copenhagen scope is stated explicitly: the formalism \emph{supports} a
Copenhagen-compatible reading; it does not prove Copenhagen is the unique
interpretation.
\end{abstract}

\tableofcontents

\section{The problem, and the choice of structure}\label{sec:choice}

\paragraph{The data.}
Fix the finite Weyl model: a system Hilbert space $\Hs$ of dimension $d^m$
($m$ registers of local dimension $d$), Weyl operators $\Wop(v)$ indexed by
$v\in V=(\mathbb Z_d)^{2m}$, with the section law
\begin{equation}\label{eq:weyl}
\Wop(u)\Wop(v)=\tau^{\,c(u,v)}\Wop(u+v),\qquad \tau=e^{i\pi/d},
\end{equation}
where $c:V\times V\to\mathbb Z_{2d}$ is the canonical $2$-cocycle of the companion
papers. A \emph{context} is a closed commuting Weyl family (an isotropic subgroup
$L\subseteq V$ with its operator family); it need not be maximal. An \emph{observer}
adjoins ancillary degrees of freedom (a friend, a memory, an apparatus) by an isometry
$\iota:\Hs\to\Hs\otimes\Hs_O$ and picks a context on the enlarged algebra.

\paragraph{Two arrow classes, one asymmetry.}
Paper D generates morphisms from two kinds of move:
\begin{enumerate}[label=(\alph*),leftmargin=2em]
\item \emph{Switches}: at a \emph{fixed} enlargement $\K$, a unitary $U$ on $\K$
carrying a context $C$ to a context $C'$ (a basis change, a Wigner rotation, the
friend performing her measurement). Switches are invertible ($U^{-1}=U^\ast$).
\item \emph{Refinements}: \emph{changing} the enlargement, $\iota\rightsquigarrow
\iota'\circ\iota$, and embedding observables by $A\mapsto A\otimes I$ (moving the
Heisenberg cut outward, from friend to Wigner). Because an isometry into a strictly
larger Hilbert space has no two-sided inverse, refinements are \emph{not} invertible.
\end{enumerate}
The single structural fact that organizes everything below is that these two classes
\emph{differ in invertibility}. Any faithful formalism must encode that asymmetry, not
merely tolerate it.

\paragraph{Three candidate structures.}
\begin{enumerate}[label=(\arabic*),leftmargin=2em]
\item \textbf{A category with a wide subgroupoid of switches.} Minimal and
unimpeachable: declare $\Oc$ a category, observe that the switches are isomorphisms
and form a wide subgroupoid, and take all holonomy invariants there. This is exactly
the corrected Paper D wording. It is \emph{true} but under-informative: it says only
that some arrows fail to be invertible, without organizing the refinement direction.
\item \textbf{A double category.} Treat switches and refinements as two 1-cell
directions (say vertical and horizontal) with $2$-cells the commuting squares
``refine-then-switch $=$ switch-then-refine.'' This records the compatibility squares
explicitly, and is the natural home for the naturality statements. But a bare double
category posits the two directions \emph{symmetrically}; the invertibility asymmetry
(switches invertible, refinements not) has to be added by hand, and one owes the
interchange law and a filling of all squares.
\item \textbf{A category (op)fibered over enlargements / a pseudofunctor.} Let the
base $\Ecat$ be the category of enlargements, let each fiber be the switch groupoid of
a fixed enlargement, and let refinements be the (op)cartesian lifts of enlargement
maps. Equivalently, a pseudofunctor $\Ctx:\Ecat\to\Gpd$, $\K\mapsto\Sw(\K)$,
$e\mapsto(-\otimes I)$.
\end{enumerate}

\paragraph{Our choice, and why.}
\emph{We adopt structure (3): $\Oc$ is the Grothendieck opfibration of a pseudofunctor
$\Ctx:\Ecat\to\Gpd$ into groupoids.} The reasons are exactly the three desiderata of
the correction:
\begin{itemize}[leftmargin=1.4em]
\item It encodes the invertibility asymmetry \emph{intrinsically}: the invertible
switches are precisely the \emph{fiber} (vertical) arrows; the non-invertible
refinements are precisely the \emph{cartesian} (horizontal) arrows over non-invertible
base maps. Invertibility is not posited; it is the statement ``fibers are groupoids.''
\item It subsumes structure (1): the wide switch subgroupoid is literally
$\bigsqcup_{\K}\Sw(\K)$, the union of the fibers, i.e. the vertical subcategory
(Theorem~\ref{thm:subgpd}).
\item It makes ``the $\U(1)$ cocycle is pulled along a functor'' literal: reindexing
$(-\otimes I)$ is a functor between fibers, and the cocycle-on-the-nose lemma
(Lemma~\ref{lem:pull}) is exactly the statement that the fiberwise cocycles form a
reindexing-compatible section (\S\ref{sec:cocycle}).
\end{itemize}
Structure (2) is not discarded: the fibration \emph{induces} a double category
(its ``framed'' / fibrant double category in the sense of Shulman
\cite{Shulman}), whose squares are the naturality squares. We use those squares in the
movable-cut diagram (\S\ref{sec:cut}) but do not take the double category as
primitive, because it would add coherence obligations (interchange) without adding
faithfulness. We record this relationship honestly in the status table
(\S\ref{sec:status}).

\begin{convention}[Strictification]\label{conv:strict}
To keep composition strictly associative we fix, once and for all, a strictly monoidal
model of enlargements: ancilla registers are indexed by a well-ordered set and
$\otimes$ is the ordered tensor, so that $(-\otimes I_{O'})\circ(-\otimes
I_{O})=(-\otimes I_{O\sqcup O'})$ \emph{on the nose}. Under this convention $\Ctx$ is a
\emph{strict} $2$-functor and the Grothendieck construction is a strict opfibration.
Dropping the convention costs only the coherence isomorphisms of a pseudofunctor, and
the passage to a fibration is then the standard Grothendieck correspondence
\cite{SGA1,JacobsBook,Johnstone}; we flag exactly where this is invoked.
\end{convention}

\section{The base: the category of enlargements \texorpdfstring{$\Ecat$}{E}}
\label{sec:base}

\begin{definition}[Enlargements]\label{def:base}
The \emph{category of enlargements} $\Ecat$ has:
\begin{itemize}[leftmargin=1.4em]
\item \textbf{Objects:} cut data, i.e. isometries $\iota:\Hs\to\Hs\otimes\Hs_{O}$ with
$\Hs$ the fixed system and $\Hs_O$ a (possibly trivial) ancilla; write the enlarged
space as $\K_\iota:=\Hs\otimes\Hs_O$. The trivial enlargement $\iota_0=\id_\Hs$ is an
object.
\item \textbf{Morphisms} $\iota\to\iota'$: ancilla-adjoining isometries
$e:\K_\iota\to\K_{\iota'}$ of the form $e=\id_{\K_\iota}\otimes\,\eta$, with
$\eta:\mathbb C\to\Hs_{O'}$ a unit vector (adjoining a fresh ancilla in a fixed state),
such that $\iota'=e\circ\iota$. Composition is composition of isometries; the identity
is $\id_{\K_\iota}$.
\end{itemize}
At the level of observables, $e$ induces the unital $\ast$-embedding
$e_\sharp:\B(\K_\iota)\to\B(\K_{\iota'})$, $A\mapsto A\otimes I$ (Heisenberg picture).
\end{definition}

\begin{proposition}[$\Ecat$ is a category; refinements are non-invertible]
\label{prop:base}
$\Ecat$ is a well-defined category. Under Convention~\ref{conv:strict} it is strictly
associative and strictly monoidal under $\otimes$. A morphism $e:\iota\to\iota'$ is
invertible in $\Ecat$ iff $\dim\K_\iota=\dim\K_{\iota'}$, i.e. iff the adjoined ancilla
is trivial. \STC
\end{proposition}
\begin{proof}
Composition of isometries is an isometry and is associative; identities act as
identities; this is a category. Strictness is Convention~\ref{conv:strict}. If
$\dim\K_{\iota'}>\dim\K_\iota$ then $e$ is a non-surjective isometry, hence has no
right inverse in $\Ecat$ (any $r$ with $e\circ r=\id$ would force $e$ surjective); if
dimensions agree, $e$ is a unitary and $e^\ast$ is its inverse. \end{proof}

The trivial enlargement $\iota_0$ is initial in the full subcategory on a fixed system
$\Hs$; the whole ``moving the cut outward'' picture is the poset-like flow away from
$\iota_0$.

\section{The fibers: switch groupoids \texorpdfstring{$\Sw(\K)$}{Sw(K)}}
\label{sec:fiber}

\begin{definition}[Switch groupoid]\label{def:sw}
For a fixed enlarged space $\K$, the \emph{switch groupoid} $\Sw(\K)$ has:
\begin{itemize}[leftmargin=1.4em]
\item \textbf{Objects:} contexts $C$ on $\B(\K)$ (closed commuting Weyl families of the
enlarged Weyl system).
\item \textbf{Morphisms} $C\to C'$: gauge classes $[U]$ of context-carrying unitaries,
i.e. $U$ in the Weyl normalizer (Clifford group) with $UCU^\ast=C'$, taken modulo the
\emph{Weyl-section gauge} $\Gamma=\{\tau^{k}\Wop(a)\}$ acting by
$U\mapsto g_2\,U\,g_1$, $g_1,g_2\in\Gamma$.
\item \textbf{Composition:} $[U_2]\circ[U_1]=[U_2U_1]$; \textbf{identity:}
$\id_C=[\,I\,]$.
\end{itemize}
\end{definition}

\begin{proposition}[$\Sw(\K)$ is a groupoid]\label{prop:fibergpd}
$\Sw(\K)$ is a well-defined groupoid: composition is associative and gauge-well-defined,
and every morphism $[U]:C\to C'$ has inverse $[U^\ast]:C'\to C$. \STC
\end{proposition}
\begin{proof}
Matrix multiplication is associative. Gauge-invariance of composition: if
$U_1\sim g_2U_1g_1$ and $U_2\sim h_2U_2h_1$, then $U_2U_1\sim
h_2U_2(h_1g_2)U_1g_1$; since $\Gamma$ is a group normalized by the Clifford group,
$h_1g_2$ is absorbed into the left/right gauge of the product, so $[U_2U_1]$ depends
only on $[U_1],[U_2]$. If $UCU^\ast=C'$ then $U^\ast C'U=C$, and
$[U^\ast][U]=[U^\ast U]=[I]=\id_C$, $[U][U^\ast]=\id_{C'}$. Hence every arrow is
invertible. \end{proof}

\begin{remark}[Why the gauge quotient]
The gauge $\Gamma$ is the phase/Weyl ambiguity of the section \eqref{eq:weyl}. Quotient\-ing
by it makes objects ``contexts'' rather than ``sectioned contexts,'' and is exactly the
freedom under which the raw holonomy $\hol$ is defined only up to sign while $\hol^2$ is
invariant (\S\ref{sec:holonomy}). The $\U(1)$ cocycle lives on this groupoid as a
central extension datum recording that quotient (\S\ref{sec:cocycle}).
\end{remark}

\section{The total category \texorpdfstring{$\Oc$}{O} as a Grothendieck opfibration}
\label{sec:total}

\begin{construction}[The reindexing pseudofunctor]\label{con:pf}
Define $\Ctx:\Ecat\to\Gpd$ on objects by $\Ctx(\iota)=\Sw(\K_\iota)$ and on a
refinement $e:\iota\to\iota'$ by the \emph{adjoin-and-tensor} functor
\[
\pushf{e}:=(-\otimes I):\Sw(\K_\iota)\to\Sw(\K_{\iota'}),\qquad
C\mapsto C\otimes I,\quad [U]\mapsto[U\otimes I].
\]
\end{construction}

\begin{lemma}[$\pushf{e}$ is a functor of groupoids]\label{lem:reindex}
$\pushf e$ is a well-defined, faithful functor. Under Convention~\ref{conv:strict},
$\pushf{(e'\circ e)}=\pushf{e'}\circ\pushf{e}$ and $\pushf{\id}=\id$, so $\Ctx$ is a
strict $2$-functor. \STC
\end{lemma}
\begin{proof}
$U\otimes I$ conjugates $C\otimes I$ to $C'\otimes I$; $(U_2U_1)\otimes
I=(U_2\otimes I)(U_1\otimes I)$ and $I\otimes I=I$, so $\pushf e$ preserves composition
and identities. It respects the gauge because $g\otimes I\in\Gamma_{\K'}$ for
$g\in\Gamma_\K$. Faithfulness: $U\otimes I=U'\otimes I$ (mod gauge) forces $[U]=[U']$.
Strict functoriality in $e$ is the definition of the ordered tensor
(Convention~\ref{conv:strict}). \end{proof}

\begin{definition}[The observer--context category]\label{def:O}
$\Oc:=\displaystyle\int_{\Ecat}\Ctx$ is the Grothendieck construction of $\Ctx$:
\begin{itemize}[leftmargin=1.4em]
\item \textbf{Objects:} pairs $o=(\iota,C)$ with $\iota\in\Ecat$ and $C\in\Sw(\K_\iota)$
a context on the enlarged algebra --- exactly Paper D's observers.
\item \textbf{Morphisms} $o=(\iota,C)\to o'=(\iota',C')$: pairs $(e,[U])$ with
$e:\iota\to\iota'$ in $\Ecat$ and $[U]:\pushf e(C)=C\otimes I\to C'$ in
$\Sw(\K_{\iota'})$.
\item \textbf{Source/target:} $\mathrm{src}(e,[U])=(\iota,C)$,
$\mathrm{tgt}(e,[U])=(\iota',C')$.
\item \textbf{Composition:} for $(e,[U]):o\to o'$ and $(e',[U']):o'\to o''$,
\begin{equation}\label{eq:comp}
(e',[U'])\circ(e,[U]):=\bigl(e'\circ e,\ [U']\circ\pushf{e'}[U]\bigr)
=\bigl(e'\circ e,\ [U'\,(U\otimes I')]\bigr).
\end{equation}
\item \textbf{Identity:} $\id_{(\iota,C)}=(\id_\iota,[I])$.
\end{itemize}
\end{definition}

\begin{theorem}[$\Oc$ is a strict opfibration]\label{thm:opfib}
Under Convention~\ref{conv:strict}, $\Oc$ is a category, the projection
$\pi:\Oc\to\Ecat$, $(\iota,C)\mapsto\iota$, $(e,[U])\mapsto e$, is a functor, and
$\pi$ is a split Grothendieck opfibration whose fiber over $\iota$ is $\Sw(\K_\iota)$.
Composition \eqref{eq:comp} is associative and unital. \STC{} (associativity, units,
opfibration split structure by direct verification); the identification with the
abstract Grothendieck construction, and the pseudofunctor version without
Convention~\ref{conv:strict}, is \STS{} \cite{SGA1,JacobsBook,Johnstone}.
\end{theorem}
\begin{proof}
\emph{Unitality.} $(\id_{\iota'},[I])\circ(e,[U])=(e,\,[I]\circ\pushf{\id}[U])=(e,[U])$
and $(e,[U])\circ(\id_\iota,[I])=(e,\,[U]\circ\pushf e[I])=(e,[U])$ since $\pushf
e[I]=[I]$.

\emph{Associativity.} Take $(e,[U]):o_0\to o_1$, $(e',[U']):o_1\to o_2$,
$(e'',[U'']):o_2\to o_3$. Composing the first two then the third gives base component
$e''\circ(e'\circ e)$ and fiber component
\[
[U'']\circ\pushf{e''}\!\bigl([U']\circ\pushf{e'}[U]\bigr)
=[U'']\circ\pushf{e''}[U']\circ\pushf{e''}\pushf{e'}[U],
\]
using functoriality of $\pushf{e''}$ (Lemma~\ref{lem:reindex}). Composing the last two
first gives base $ (e''\circ e')\circ e$ and fiber
$[U'']\circ\pushf{e''}[U']\circ\pushf{(e''\circ e')}[U]$. By Convention~\ref{conv:strict},
$e''\circ(e'\circ e)=(e''\circ e')\circ e$ and
$\pushf{e''}\pushf{e'}=\pushf{(e''\circ e')}$, so the two agree. Hence composition is
associative.

\emph{Functoriality and fibers.} $\pi$ evidently preserves composition and identities,
and $\pi^{-1}(\iota)$ (arrows over $\id_\iota$) is $\Sw(\K_\iota)$.

\emph{Opcartesian lifts.} For a base map $e:\iota\to\iota'$ and object
$(\iota,C)$, the arrow $\rho_e:=(e,[I]):(\iota,C)\to(\iota',\pushf e C)$ is opcartesian:
any $(e'\circ e,[V]):(\iota,C)\to(\iota'',C'')$ factors uniquely through $\rho_e$ as
$(e',[V])\circ\rho_e$. The assignment $e\mapsto\rho_e$ is a cleavage, split because
$\pushf{}$ is strictly functorial. This is the Grothendieck construction of a strict
$2$-functor \cite{SGA1,JacobsBook}. \end{proof}

\begin{proposition}[Opcartesian--vertical factorization]\label{prop:fact}
Every morphism of $\Oc$ factors uniquely as a refinement followed by a switch:
\[
(e,[U])=\underbrace{(\id_{\iota'},[U])}_{\text{switch}}\ \circ\
\underbrace{(e,[I])}_{\text{refinement}}.
\]
Refinements are the opcartesian arrows $\rho_e=(e,[I])$; switches are the vertical
arrows $(\id_\iota,[U])$. \STC
\end{proposition}
\begin{proof}
$(\id_{\iota'},[U])\circ(e,[I])=(e,\,[U]\circ\pushf e[I])=(e,[U])$. Uniqueness of the
(opcartesian, vertical) factorization is the defining property of an opfibration
(Theorem~\ref{thm:opfib}). \end{proof}

\subsection*{Dictionary}
\begin{center}\small
\begin{tabular}{lll}
\toprule
Paper D term & Structure in $\Oc$ & Invertible?\\
\midrule
observer $(\iota,C)$ & object of $\Oc$ / object of fiber $\Sw(\K_\iota)$ & ---\\
switch $C\to C'$ (fixed cut) & vertical arrow $(\id_\iota,[U])$ & yes (fiber = groupoid)\\
refinement $A\mapsto A\otimes I$ & opcartesian arrow $(e,[I])$ & no (Prop.~\ref{prop:base})\\
``moving the cut'' & travel along the base $\Ecat$ & no\\
switch loop & loop in a single fiber $\Sw(\K)$ & yes (holonomy defined)\\
\bottomrule
\end{tabular}
\end{center}

\section{Switches form a wide subgroupoid}\label{sec:subgpd}

\begin{theorem}[Wide switch subgroupoid]\label{thm:subgpd}
The vertical arrows $\Sw(\Oc):=\{(\id_\iota,[U])\}$ form a wide subgroupoid of $\Oc$:
it contains every object, is closed under composition and identities, and every one of
its arrows is invertible, with $(\id_\iota,[U])^{-1}=(\id_\iota,[U^\ast])$. As a
category $\Sw(\Oc)=\bigsqcup_{\iota\in\Ecat}\Sw(\K_\iota)$, the disjoint union of the
fibers. \STC
\end{theorem}
\begin{proof}
Verticality is preserved by \eqref{eq:comp}: if both base components are identities so
is their composite, and $\pushf{\id}=\id$ reduces the fiber component to $[U']\circ[U]$,
composition inside $\Sw(\K_\iota)$. Identities $(\id_\iota,[I])$ are vertical.
Invertibility is Proposition~\ref{prop:fibergpd} transported to $\Oc$:
$(\id_\iota,[U^\ast])\circ(\id_\iota,[U])=(\id_\iota,[U^\ast U])=(\id_\iota,[I])$.
Wideness holds because every object $(\iota,C)$ is in the fiber $\Sw(\K_\iota)$. There
are no vertical arrows between different fibers (their base components would differ from
$\id$), so the subcategory is the disjoint union of fibers. \end{proof}

\noindent This is precisely candidate structure (1) of \S\ref{sec:choice}, now
\emph{derived} rather than posited: ``category with a wide switch subgroupoid'' is the
underlying category plus its vertical subcategory. All holonomy invariants of
\S\ref{sec:holonomy} live here.

\section{Refinement functoriality}\label{sec:refine}

Fix a refinement $e:\iota\to\iota'$ with reindexing $\pushf e=(-\otimes I)$. We record
its action on each layer of structure.

\paragraph{(R1) On contexts and switches. \STC}
By Lemma~\ref{lem:reindex}, $\pushf e:\Sw(\K_\iota)\to\Sw(\K_{\iota'})$ is a faithful
functor of groupoids: $C\mapsto C\otimes I$, $[U]\mapsto[U\otimes I]$. It is injective
on objects and faithful on arrows, but not full (the enlarged algebra has contexts and
switches touching the ancilla that are not of the form $-\otimes I$). Thus refinement
embeds each observer's switch groupoid into the larger one.

\paragraph{(R2) On the $\U(1)$ cocycle (on the nose). \STC}
Let $c_{\K}$ denote the enlarged Weyl cocycle \eqref{eq:weyl} on $\Sw(\K)$.

\begin{lemma}[Pullback of the cocycle]\label{lem:pull}
$\pushf e$ pulls the cocycle back strictly: $\pushf e^{\,\ast}c_{\K_{\iota'}}=c_{\K_\iota}$.
Equivalently $(\Wop(u)\otimes I)(\Wop(v)\otimes I)=\tau^{c(u,v)}(\Wop(u+v)\otimes I)$
with the \emph{same} exponent $c(u,v)$, for any enlargement. \STC
\end{lemma}
\begin{proof}
$(\Wop(u)\otimes I)(\Wop(v)\otimes I)=\Wop(u)\Wop(v)\otimes
I=\tau^{c(u,v)}\Wop(u+v)\otimes I$ by \eqref{eq:weyl}; the phase $\tau^{c(u,v)}$ is
unchanged, so the cocycle is preserved verbatim. Machine certificate: the six embedded
Peres--Mermin context products at $d=8$ are $(+1)^{\times5},(-1)^{\times1}$, identical
to $d=4$ (\texttt{observer\_groupoid.py}, part~A). \end{proof}

\paragraph{(R3) On value assignments (pullback--monotonicity). \STC}
A \emph{value assignment} on a family $F$ of contexts is a section of the presheaf of
$\mathbb Z_d$-valued (or $\U(1)$-valued) noncontextual valuations.

\begin{lemma}[Restriction and monotonicity]\label{lem:mono}
Refinement induces a restriction map $\pushf e^{\,\ast}$ on value assignments: a value
assignment on $\pushf e(F)=F\otimes I$ restricts along $\pushf e$ to one on $F$. Hence
\[
\text{$F$ has no global value assignment}\ \Longrightarrow\
\text{no enlarged family containing $\pushf e(F)$ has one.}
\]
In particular the Peres--Mermin obstruction survives every Wigner enlargement. \STC
\end{lemma}
\begin{proof}
$\pushf e$ is an injective strict morphism of context structures preserving all
operator products (Lemma~\ref{lem:pull}); a valuation $\nu$ on a family containing
$\pushf e(F)$ pulls back to $\nu\circ\pushf e$ on $F$, which is again noncontextual and
consistent because products are preserved on the nose. Contrapositive gives the
monotonicity. Certificate as in Lemma~\ref{lem:pull}. \end{proof}

\paragraph{(R4) On obstruction classes (pullback--monotonicity of the class). \STC/\STS}
Let $[\theta_F]\in H^2(-;\U(1))$ (equivalently the $\mathbb Z_d$/AvN obstruction of the
companion papers) be the obstruction class of a family.

\begin{proposition}[Class pullback]\label{prop:classpull}
$\pushf e^{\,\ast}[\theta_{\pushf e(F)}]=[\theta_F]$: the obstruction class of the
embedded family equals the pullback of the class upstairs. Consequently the class is
\emph{monotone} under refinement --- it cannot be trivialized by enlarging the cut ---
and $\pushf e$ never creates or destroys the AvN sector. \STC{} for the equality of
representing cocycles (Lemma~\ref{lem:pull}); \STS{} for the statement at the level of
cohomology classes, which uses functoriality of $H^2(-;\U(1))$ under the groupoid
functor $\pushf e$ \cite[naturality of group/groupoid cohomology]{Brown,Weibel}.
\end{proposition}
\begin{proof}
Cocycle representatives are preserved by Lemma~\ref{lem:pull}; naturality of $H^2$
under $\pushf e$ turns this into an equality of classes. Monotonicity is the
contrapositive together with (R3). \end{proof}

\begin{remark}[Direction of variance]
Refinement $\pushf e$ is \emph{covariant} on contexts/switches ($C\mapsto C\otimes I$)
but the induced map on cocycles/classes is \emph{contravariant} (pullback). This is why
the natural home is an \emph{op}fibration (covariant reindexing of fibers) carrying a
\emph{contravariant} cocycle section --- see \S\ref{sec:cocycle}.
\end{remark}

\section{Where the \texorpdfstring{$\U(1)$}{U(1)} cocycle lives}\label{sec:cocycle}

We can now state precisely where the cocycle lives and why.

\begin{definition}[Fiberwise cocycle / reindexing-compatible section]\label{def:cocycle}
For each enlargement $\iota$ let $c_{\K_\iota}\in Z^2(\Sw(\K_\iota);\U(1))$ be the
enlarged Weyl cocycle, i.e. the $2$-cocycle of the central extension
$\U(1)\to\U(\K_\iota)\to\PU(\K_\iota)$ restricted along the switch groupoid, normalized
on Weyl operators by \eqref{eq:weyl}. The family
$\{c_{\K_\iota}\}_{\iota\in\Ecat}$ is a \emph{reindexing-compatible section} of the
contravariant assignment $\iota\mapsto Z^2(\Sw(\K_\iota);\U(1))$: by
Lemma~\ref{lem:pull},
\begin{equation}\label{eq:section}
\pushf e^{\,\ast}c_{\K_{\iota'}}=c_{\K_\iota}\qquad\text{for every refinement } e.
\end{equation}
\end{definition}

\begin{theorem}[Location of the cocycle]\label{thm:location}
The $\U(1)$ cocycle lives \emph{on the switch subgroupoid, fiberwise, as the
reindexing-compatible section \eqref{eq:section}} --- equivalently, as a $\U(1)$-central
extension of the vertical subcategory $\Sw(\Oc)$ that is refinement-equivariant. It does
\emph{not} live on the whole category $\Oc$ as an ordinary group/groupoid $2$-cocycle.
\STC{} (fiberwise existence and compatibility). The assembly into a single class on all
of $\Oc$ is \STG{} (see below).
\end{theorem}
\begin{proof}[Justification]
\emph{Why the switch subgroupoid.} A $\U(1)$ $2$-cocycle with holonomy is the datum of
a central extension, and holonomy of a loop requires the loop to be \emph{invertible}.
Only the vertical (switch) arrows are invertible (Theorem~\ref{thm:subgpd}); the
non-invertible refinements cannot carry holonomy, only pull the cocycle back
(Lemma~\ref{lem:pull}). Hence the natural carrier is $\Sw(\Oc)=\bigsqcup_\iota
\Sw(\K_\iota)$, one central extension per fiber.

\emph{Why fiberwise-plus-compatible rather than global.} A single $2$-cocycle on the
whole category $\Oc$ would be an element of a cohomology theory of $\Oc$ that accepts
non-invertible arrows --- Baues--Wirsching or Thomason cohomology of the category, i.e.
$H^2$ of the nerve $N\Oc$ with the appropriate coefficient system \cite{BauesWirsching,
Thomason}. Ordinary groupoid $2$-cocycles are not defined on non-invertible arrows.
What we \emph{have} constructed is the next best and arguably correct object: a cocycle
on each fiber, glued by the reindexing pullback \eqref{eq:section}. This is exactly a
section of the fibered cohomology of $\pi$, and it is what ``the cocycle is pulled along
a functor'' means. \end{proof}

\begin{remark}[The gerbe / nerve viewpoint --- \STG]\label{rem:gerbe}
Two stronger packagings are natural but not established here. (a) The compatible section
\eqref{eq:section} is the cocycle data of a \emph{$\U(1)$-gerbe} over the base $\Ecat$
(a stack of $\U(1)$-central extensions of the switch groupoids); proving it satisfies
the gerbe descent/coherence conditions is a \STG. (b) A single class in
$H^2(N\Oc;\U(1))$ (Baues--Wirsching) restricting to $c_{\K_\iota}$ on each fiber would
be the ideal ``one cohomological object'' the program advertises; constructing it, or
proving the obstruction to its existence, is the \textbf{main open architectural
question} (\S\ref{sec:status}). Everything downstream (holonomy, movable cut) uses only
the fiberwise section, which is complete.
\end{remark}

\section{Holonomy, and its distinction from pullback}\label{sec:holonomy}

\begin{definition}[Switch-loop holonomy]\label{def:hol}
Let $\gamma=[U_k]\circ\cdots\circ[U_1]$ be a loop in a single fiber $\Sw(\K)$ (so
$U_k\cdots U_1$ carries a context to itself). Choose a special-unitary closer $W_c$ and
set
\[
\hol(\gamma)=\det\!\bigl(W_c\,U_k\cdots U_1\bigr),\qquad
\hol^2(\gamma)=\det\!\bigl(U_k\cdots U_1\bigr)^2 .
\]
\end{definition}

\begin{lemma}[Holonomy is a fiber invariant; $\hol^2$ is gauge-robust]\label{lem:holinv}
$\hol$ is defined for invertible switch loops only. $\hol$ itself is section-dependent
(changes by $\pm1$ under Weyl regauging of any switch and by the choice of closer), but
$\hol^2\in\{\pm1\}$ is invariant under Weyl regauging and closer choice, hence a
well-defined $\Zt$-valued function on switch loops. \STC{} (Paper D
Lemma~2; every Weyl element has $\det=\pm1$ and closers are special unitary;
machine-checked over all $16$ Weyl elements, \texttt{observer\_groupoid.py} part~B).
\end{lemma}

\begin{proposition}[Holonomy vs.\ pullback are different operations]\label{prop:holvspull}
Let $\gamma$ be a switch loop in $\Sw(\K_\iota)$ and $e:\iota\to\iota'$ a refinement.
Then:
\begin{enumerate}[label=(\roman*),leftmargin=2em]
\item \emph{Holonomy} is the trace of going \emph{around an invertible loop inside one
fiber}; it is a number in $\U(1)$ (resp. $\hol^2\in\Zt$).
\item \emph{Pullback along a refinement} is the functor $\pushf e$ \emph{between
different fibers}; a refinement is not a loop and has no holonomy of its own.
\item The two interact by \emph{preservation}: $\hol(\pushf e\gamma)=\hol(\gamma)$ and
$\hol^2(\pushf e\gamma)=\hol^2(\gamma)$, because $\pushf e$ preserves the cocycle on the
nose (Lemma~\ref{lem:pull}) and $\det(A\otimes I_n)=\det(A)^n$ with the closer
scaled accordingly, so the $\Zt$ residue is unchanged.
\end{enumerate}
\STC{} (i)--(ii) by definition; (iii) \STC{} for the $\hol^2$ residue, \STS{} for the
raw $\hol$ phase where the ancilla dimension enters $\det(A\otimes I_n)=\det(A)^n$ and
one must track the closer normalization.
\end{proposition}

\noindent
The upshot, which the ``groupoid'' error obscured: \emph{moving the cut is not
traversing a loop.} Refinement is pullback-along-an-embedding between fibers; it
transports holonomy but is not itself a holonomy. Only genuine switch loops --- the
friend rotating her basis and returning --- have holonomy. This is why holonomy is
defined solely on the switch subgroupoid.

\section{The movable cut, exactly}\label{sec:cut}

We restate Paper D's Theorem~3 as an exact multi-level statement. Fix an observer
$o=(\iota,C)$ (the \emph{friend} description), a refinement $e:\iota\to\iota'$ to
$o^\uparrow=(\iota',C\otimes I)$ (the \emph{Wigner} description), and a switch
$s=[U]:C\to C'$ in the friend fiber (the friend's basis change / measurement).
Refinement transports $s$ to $\pushf e s=[U\otimes I]:C\otimes I\to C'\otimes I$.

\begin{theorem}[Movable cut, five separated levels]\label{thm:cut}
With the data above, and any ancilla state $\sigma$ ($\tr\sigma=1$):
\begin{enumerate}[label=\textbf{(\arabic*)},leftmargin=2.6em]
\item \textbf{Empirical-probability invariance \STC.} For every state $\rho$ on
$\K_\iota$ and every outcome effect $E$ measured in the switched context,
\[
\tr(\rho\,E)=\tr\!\bigl((\rho\otimes\sigma)\,(E\otimes I)\bigr).
\]
Refinement changes no Born probability; friend and Wigner assign identical statistics.
\item \textbf{Ray-level (projective) equivalence \STC.} The square below commutes
\emph{strictly} in $\Oc$ (equivalently at the $\PU$/density-operator level): the
comparison map $R=\rho_e$ intertwines the friend and Wigner descriptions, and the
ray-level holonomy of the switch loop is $+1$ (hardware-measured on \texttt{ibm\_fez}:
loop phase $-2.9^\circ\pm1.6^\circ$).
\item \textbf{Section-dependent phase \STC.} Lifting rays to vectors (choosing a
determinant section) makes the comparison defined only up to $\U(1)$: around a closed
switch loop $\gamma$ the section phase is exactly the holonomy $\hol(\gamma)$
(Definition~\ref{def:hol}); for the mutually-unbiased triangle $\hol=-i$ in the Clifford
section. This phase is bookkeeping, not statistics.
\item \textbf{Cohomology-class invariance \STC.} The gauge-robust residue
$\hol^2(\gamma)\in\Zt$ and the obstruction class $[\theta]$ are unchanged by the move
and by re-sectioning (Lemma~\ref{lem:holinv}, Proposition~\ref{prop:classpull}).
\item \textbf{Interpretive conclusion --- \emph{interpretation, not theorem}.} ``The
cut may be moved without empirical consequence, the section bookkeeping recording the
move as a $\U(1)$ phase whose invariant residue is the class.'' This is the reading the
formalism \emph{supports}; it is labeled interpretation and is not claimed unique
(\S\ref{sec:copenhagen}).
\end{enumerate}
\end{theorem}

\begin{proof}
\textbf{(1)} $\tr((\rho\otimes\sigma)(E\otimes I))=\tr(\rho E)\,\tr(\sigma)=\tr(\rho E)$.
\textbf{(2)} Strict commutativity of the square is
Proposition~\ref{prop:fact}/\eqref{eq:comp}: $(\id_{\iota'},[U\otimes
I])\circ(e,[I])=(e,[U\otimes I])=(e,\pushf e[U])=\rho_e\circ(\text{lift of }s)$ at the
level of gauge classes (rays). The measured ray-level loop phase is $+1$ by
Lemma~\ref{lem:holinv} applied at the $\PU$ level.
\textbf{(3)} Choosing a determinant section is a set-theoretic lift
$\PU\to\U$; its failure to be a functor around $\gamma$ is precisely $\hol(\gamma)$
(Definition~\ref{def:hol}), section-dependent by Lemma~\ref{lem:holinv}.
\textbf{(4)} $\hol^2$ and $[\theta]$ are the gauge-/section-invariants of
Lemma~\ref{lem:holinv} and Proposition~\ref{prop:classpull}.
\textbf{(5)} is flagged as interpretation. \end{proof}

\begin{center}
\begin{tikzcd}[row sep=large, column sep=huge]
(\iota,\,C)
  \arrow[r, "{s=[U]}"]
  \arrow[d, "{R=\rho_e}"']
& (\iota,\,C')
  \arrow[d, "{R=\rho_e}"]
\\
(\iota',\,C\otimes I)
  \arrow[r, "{\pushf{e}\,s=[U\otimes I]}"']
& (\iota',\,C'\otimes I)
\end{tikzcd}
\end{center}
\noindent\emph{Reading of the diagram.} Horizontal arrows are switches (friend's move,
invertible); vertical arrows are the refinement (moving the cut, non-invertible).
\textbf{At ray level} (Theorem~\ref{thm:cut}(2)) the square commutes on the nose.
\textbf{At vector/section level} (Theorem~\ref{thm:cut}(3)) it commutes up to a $\U(1)$
phase; \textbf{around a closed switch loop} that phase is $\hol$, with invariant residue
$\hol^2$ (Theorem~\ref{thm:cut}(4)). The square is the naturality square of the induced
double category (Shulman \cite{Shulman}); we use it here only as a picture of the exact
statement.

\begin{remark}[What is \emph{not} proved]
Theorem~\ref{thm:cut} proves invariance of statistics, of the class, and the exact
section-phase bookkeeping for \emph{this} friend/Wigner pair and \emph{these} moves. It
does \emph{not} prove the fully general Copenhagen ``movable cut principle'' for
arbitrary system/apparatus splits, nor that the update rule is unique beyond the
constrained affine depolarizing family of the companion note (whose unique single-Kraus
member is L\"uders). Those are the boundaries of the claim.
\end{remark}

\section{Copenhagen scope}\label{sec:copenhagen}

\begin{proposition}[Scope, explicit]\label{prop:scope}
The architecture of \S\S\ref{sec:choice}--\ref{sec:cut} is
\emph{Copenhagen-compatible}: local classicality per context, complementarity as
holonomy of switch loops, relative facts as fiberwise value assignments obstructed
globally, and the movable cut as pullback-along-refinement, are all theorems of the
stated finite-Weyl axioms. It does \emph{not} establish that Copenhagen is the unique
interpretation consistent with these structures. Everything proved here is available to
rival interpretations verbatim; what is fixed is the mathematics (an opfibration with a
fiberwise $\U(1)$ cocycle), not the metaphysics. \STC{} (as a scope statement).
\end{proposition}

Concretely: nothing above modifies a Born-rule prediction; the constructions are
reorganizations of finite quantum data. A many-worlds or Bohmian reader may accept
Theorems~\ref{thm:subgpd}--\ref{thm:cut} without accepting the interpretive gloss
\ref{thm:cut}(5). We therefore present the results as \emph{structural
characterizations within the axioms}, not as an interpretation-free derivation of
Copenhagen.

\section{Status, and the shape of the object}\label{sec:status}

\subsection*{Status table}
\begin{center}\small
\begin{tabular}{p{7.9cm}l}
\toprule
Ingredient & Status\\
\midrule
Base category $\Ecat$ of enlargements; non-invertibility of refinements (Prop.~\ref{prop:base}) & \STC\\
Fibers $\Sw(\K)$ are groupoids (Prop.~\ref{prop:fibergpd}) & \STC\\
Reindexing $\pushf e=(-\otimes I)$ is a faithful functor; $\Ctx$ strict $2$-functor (Lem.~\ref{lem:reindex}) & \STC\\
Objects, arrows, src/tgt, composition, identities (Def.~\ref{def:O}) & \STC\\
Associativity and unitality of $\Oc$ (Thm.~\ref{thm:opfib}) & \STC\\
$\Oc=\int_\Ecat\Ctx$ is a split opfibration (Thm.~\ref{thm:opfib}) & \STS{} \cite{SGA1,JacobsBook,Johnstone}\\
Pseudofunctor version without strictification (Conv.~\ref{conv:strict}) & \STS{} \cite{SGA1,Johnstone}\\
Switches form a wide subgroupoid $=\bigsqcup_\iota\Sw(\K_\iota)$ (Thm.~\ref{thm:subgpd}) & \STC\\
Opcartesian--vertical factorization (Prop.~\ref{prop:fact}) & \STC\\
Refinement action on contexts/switches (R1) & \STC\\
Cocycle pulled back on the nose (Lem.~\ref{lem:pull}; certificate $d=8$) & \STC\\
Value-assignment restriction \& monotonicity (Lem.~\ref{lem:mono}) & \STC\\
Obstruction-\emph{class} pullback (Prop.~\ref{prop:classpull}) & \STC{} cocycle / \STS{} class \cite{Brown,Weibel}\\
Cocycle located: fiberwise, switch subgroupoid, compatible section (Thm.~\ref{thm:location}) & \STC\\
Single class on all of $\Oc$ ($H^2(N\Oc;\U(1))$, Baues--Wirsching) (Rem.~\ref{rem:gerbe}) & \STG\\
$\U(1)$-gerbe over $\Ecat$ packaging (Rem.~\ref{rem:gerbe}) & \STG\\
Holonomy defined on switch loops; $\hol^2\in\Zt$ gauge-robust (Lem.~\ref{lem:holinv}) & \STC{} \cite{paperD}\\
Holonomy $\neq$ pullback; preservation $\hol(\pushf e\gamma)=\hol(\gamma)$ (Prop.~\ref{prop:holvspull}) & \STC{}/\STS\\
Movable cut, levels (1)--(4) (Thm.~\ref{thm:cut}) & \STC\\
Movable cut, level (5) interpretive & interpretation\\
Copenhagen scope statement (Prop.~\ref{prop:scope}) & \STC\\
Induced (framed) double category of squares (Shulman) & \STS{} \cite{Shulman}\\
\bottomrule
\end{tabular}
\end{center}

\subsection*{Is it a category, double category, bicategory, or fibration?}

\paragraph{Best judgment.}
The correct minimal-yet-faithful structure is a \textbf{split Grothendieck opfibration
$\pi:\Oc\to\Ecat$ over the category of enlargements, equivalently a (strict, after
Convention~\ref{conv:strict}) pseudofunctor $\Ctx:\Ecat\to\Gpd$ into groupoids},
carrying a \emph{fiberwise} $\U(1)$ cocycle compatible under reindexing
\eqref{eq:section}. Concretely:
\begin{itemize}[leftmargin=1.4em]
\item It is \emph{not merely} a category: calling it ``a category with a wide switch
subgroupoid'' (candidate 1) is \emph{true} (Theorem~\ref{thm:subgpd}) but forgets the
base direction that makes the refinements organized rather than ad hoc.
\item It is \emph{not primitively} a double category or bicategory: the double category
(candidate 2, Shulman's framed double category of the fibration) is \emph{induced}, and
its squares are the movable-cut naturality squares, but taking it as primitive imposes
interchange coherence that the fibration gives for free and hides the invertibility
asymmetry.
\item It \emph{is} a fibration in the honest sense that the invertible/non-invertible
split is exactly the vertical/opcartesian split, and the cocycle's ``pulled along a
functor'' behavior is exactly reindexing.
\end{itemize}

\paragraph{The main open architectural question.}
Whether the fiberwise, reindexing-compatible cocycle \eqref{eq:section} \emph{assembles
into a single cohomology class on all of $\Oc$}. Precisely: does there exist a class in
the Baues--Wirsching / Thomason cohomology $H^2(N\Oc;\U(1))$ (a cohomology theory
admitting the non-invertible refinement arrows) restricting to the canonical Weyl
cocycle $c_{\K_\iota}$ on every fiber, and if so is it unique; equivalently, is the
compatible section \eqref{eq:section} the cocycle datum of a genuine $\U(1)$-gerbe over
$\Ecat$? Our results need only the fiberwise section, which is complete; the
single-object global class is the natural next target and is currently a \STG. A
secondary open question is whether the raw phase $\hol$ (not merely $\hol^2$) admits a
section-free definition once the ancilla-dimension bookkeeping of
Proposition~\ref{prop:holvspull}(iii) is normalized functorially.

\begin{thebibliography}{99}\small
\bibitem{paperD} M.\ Flores Gordillo, \emph{The Observer--Context Category: Relative
Facts, a $\mathbb Z_2$ Switching Invariant, and the Movable Cut}, companion Paper~D,
2026; verification suite \texttt{observer\_groupoid.py},
\texttt{github.com/manuflog/contextuality-obstructions}.
\bibitem{note} M.\ Flores Gordillo, \emph{Local Validity and Contextual Holonomy: A
Structural Formulation of the Copenhagen Principle} (note v3), 2026; companion
Papers~A--C.
\bibitem{SGA1} A.\ Grothendieck, \emph{Rev\^etements \'etales et groupe fondamental
(SGA 1)}, Lecture Notes in Math.\ \textbf{224}, Springer, 1971 (Grothendieck
construction; fibered categories).
\bibitem{JacobsBook} B.\ Jacobs, \emph{Categorical Logic and Type Theory}, Studies in
Logic \textbf{141}, Elsevier, 1999 (fibrations, cleavages, split (op)fibrations).
\bibitem{Johnstone} P.\ T.\ Johnstone, \emph{Sketches of an Elephant: A Topos Theory
Compendium}, Vol.~1, Oxford, 2002 (fibrations $\leftrightarrow$ pseudofunctors).
\bibitem{Shulman} M.\ Shulman, \emph{Framed bicategories and monoidal fibrations},
Theory Appl.\ Categ.\ \textbf{20} (2008), 650--738 (double category induced by a
fibration).
\bibitem{BauesWirsching} H.-J.\ Baues, G.\ Wirsching, \emph{Cohomology of small
categories}, J.\ Pure Appl.\ Algebra \textbf{38} (1985), 187--211.
\bibitem{Thomason} R.\ W.\ Thomason, \emph{Homotopy colimits in the category of small
categories}, Math.\ Proc.\ Cambridge Philos.\ Soc.\ \textbf{85} (1979), 91--109.
\bibitem{Brown} K.\ S.\ Brown, \emph{Cohomology of Groups}, GTM \textbf{87}, Springer,
1982 (functoriality of group/groupoid cohomology).
\bibitem{Weibel} C.\ A.\ Weibel, \emph{An Introduction to Homological Algebra},
Cambridge, 1994.
\bibitem{KS} S.\ Kochen, E.\ Specker, \emph{The problem of hidden variables in quantum
mechanics}, J.\ Math.\ Mech.\ \textbf{17} (1967), 59--87.
\bibitem{Wigner} E.\ P.\ Wigner, \emph{Remarks on the mind--body question}, in
\emph{The Scientist Speculates} (1961).
\end{thebibliography}

\end{document}
